\subsubsection{Entanglement vs. Quantum Correlation}\label{sec:entanglement-vs-quantum-correlation}

Throughout the discussion of the EPR paradox, we have repeatedly encountered two closely related concepts: \textbf{entanglement} and \textbf{quantum correlation}. Although they are strongly connected, they describe different aspects of a quantum system.

\highspace
\textcolor{Green3}{\faIcon{link} \textbf{Entanglement.}} We have already defined \textbf{entanglement} (\autopageref{sec:entanglement}) as a property of the quantum state itself. A system is entangled when its state cannot be written as a tensor product of the states of its individual subsystems.

\highspace
For example, the Bell state:
\begin{equation*}
    \ket{\Phi^+}
    =
    \frac{\ket{00}+\ket{11}}{\sqrt{2}}
\end{equation*}
Is entangled because it cannot be expressed as:
\begin{equation*}
    \ket{\psi_A}\otimes\ket{\psi_B}
\end{equation*}
The important point is that entanglement is an intrinsic property of the system. It exists independently of whether anybody performs a measurement (see the discussion on \autopageref{question:entanglement-intrinsic-or-basis-dependent}).

\highspace
\textcolor{Green3}{\faIcon{chart-line} \textbf{Quantum Correlation.}} \definition{Quantum Correlation} refers to the \hl{statistical relationship observed between measurement outcomes}. Unlike entanglement, quantum correlation is \textbf{not a property of the state} alone. It \textbf{depends on how the system is measured}. For instance, Alice and Bob may share the same entangled Bell pair, yet the observed correlations change depending on the measurement bases they choose.

\highspace
\begin{flushleft}
    \textcolor{DarkOrange3}{\faIcon{balance-scale} \textbf{Entanglement vs. Quantum Correlation}}
\end{flushleft}
\hl{Entanglement} does not depend on the measurement basis:
\begin{itemize}
    \item Entanglement is a \hl{property of the quantum system}.
    \item It exists before any measurement is performed.
\end{itemize}
\hl{Quantum correlation}, on the other hand, depends on the measurement basis:
\begin{itemize}
    \item Quantum correlation is a \hl{property of the measurement outcomes}.
    \item Different measurement bases may produce different observed correlations.
\end{itemize}
In summary:
\begin{itemize}
    \item \important{Entanglement} is a property of the quantum state and does not depend on how the system is measured.
    \item \important{Quantum correlation} is a property of the measurement outcomes and depends on the measurement basis chosen by the observers.
\end{itemize}